\section{Scope, Method, and Qualifications}

\subsection*{Subject, Passage Boundary, and Purpose}

Genesis 15:5, 12, and 17--18 are tested here against one concrete exegetical claim: whether the chapter's narrated sequence places God's command to look at the stars in verse 5 under daylight before the sunset reported in verses 12 and 17. The inquiry examines the Masoretic Hebrew text, the Greek Septuagint, the Latin Vulgate, the Syriac Peshitta, Second Temple retellings in \emph{Jubilees} and Josephus, Don Isaac Abarbanel's Hebrew prefatory Questions 7 and 16, patristic and Catholic Latin exegesis, and the New Testament echoes of Genesis 15:5--6. The closing section offers a Catholic spiritual reflection explicitly distinguished from philological and historical demonstration.

\subsection*{Editions, Translations, and Direct-Collation Boundaries}

All English Scripture quotations are reproduced verbatim from the public-domain Challoner revision of the Douay-Rheims Bible, collated with the Clementine Vulgate. The Greek Septuagint is verified directly from Henry Barclay Swete's Cambridge edition (vol.~I, 4th ed., 1909, following Codex Alexandrinus across Codex Vaticanus's opening lacuna Gen 1:1--46:28a), set alongside Lancelot Charles Lee Brenton's English translation (1844) of the Sixtine text. Masoretic Hebrew verbal morphology for Genesis 15:5, 12, and 17--18 rests on Rudolf Kittel and Georg Beer's \emph{Biblia Hebraica}, Pars I (Leipzig, 1909), pp.~18--19, with the morphological ambiguity of \latin{b\={a}'\={a}h} in 15:17 (penultimate-stressed Qal perfect versus ultimate-stressed feminine singular active participle) stated transparently alongside both the Kittel record's participial analysis and classical Masoretic accentuation.

Every commentator analyzed in Sections 3 through 6 was inspected directly in complete or passage-checked original-language or public-domain editions: Isaac Abarbanel's Hebrew opening formula and prefatory Questions 7 and 16 on Genesis 15 (Warsaw, 1862; the verified excerpt does not contain Abarbanel's body answers, and earlier Jewish traditions---the visionary school, \emph{Shabbat} 156a, and the Ran---are reported solely as Abarbanel himself cites them within those two Hebrew questions); Cornelius a Lapide's Latin notes on Genesis 15:1, 5, 11, 12, and 17 (Antwerp, 1700); St. Augustine's \emph{De Civitate Dei} X.8 and XVI.23--24 (CSEL 40/2 Latin and Marcus Dods's 1871 English) and \emph{De Genesi ad litteram} XII.6 (PL 34); St. Jerome's \emph{Quaestiones Hebraicae in Genesim} ch.~15 (PL 23); St. Bede's \emph{In principium Genesis} on ch.~15 (PL 91); Alcuin of York's \emph{Interrogationes} QQ.~165--170 (PL 100); Angelomus of Luxeuil's \emph{Commentarius in Genesim} ch.~15 (PL 115); St. Clement of Rome's \emph{1 Clement} 10 (ANF 1); St. Ambrose's \emph{De excessu fratris Satyri} II.95 (NPNF2 10); St. John Chrysostom's \emph{Homilies on Romans} 9 (NPNF1 11) and \emph{Commentary on Galatians} on 3:6 (NPNF1 13); and St. Thomas Aquinas's \emph{Summa Theologiae} II-II q.~173 a.~3, q.~174 a.~6, III q.~30 a.~3 ad~3 and \emph{Super Galatas} 3:16. The Syriac Peshitta (15:5, 12, 17--18), \emph{Book of Jubilees} 14:1--20 (Charles, 1913), and Josephus's \emph{Antiquitates Judaicae} I.10.3 (Whiston, 1737) were verified through inspected passage records.

Other ancient and medieval works bearing on Genesis 15---including the Aramaic Targums (Onkelos and Pseudo-Jonathan), Philo's \emph{Quis rerum divinarum heres sit}, Babylonian Talmud \emph{Shabbat} 156a and \emph{Pesachim} 6b in their own tractate texts, \emph{Bereshit Rabbah} 44:12, medieval and modern Jewish commentaries other than Abarbanel's Questions 7 and 16 (Rashi, Ibn Ezra, Maimonides, Radak, Ramban, Chizkuni, Rabbeinu Bahya, Jacob ben Asher, Malbim, and the Netziv), George Leo Haydock's Bible notes, and the uncollated patristic Genesis commentaries of St. John Chrysostom (\emph{Homilies on Genesis} 36--37), St. Ephrem the Syrian, Theodoret of Cyrus, St. Ambrose (\emph{De Abraham}), and St. Augustine (\emph{Retractationes} II.43)---are not held locally as directly inspectable Genesis 15 texts and are therefore excluded from body analysis and claims of verified silence.

\subsection*{Transliteration, Rights, and Theological Qualification}

Hebrew and Greek words appear in Latin-script transliteration (\latin{\={a}, \={e}, \={\i}, \={o}, \={u}, \d{h}, \d{s}, \d{t}}) with inline English renderings. All verbatim quotations come from public-domain texts or direct translations from public-domain Latin and Hebrew editions; no noncommercial or copyrighted modern Scripture or rabbinic translation is reproduced.

Offered as an aid to Catholic scriptural study and prayerful reflection, the text is not a magisterial act or a critical edition. Where Catholic doctrine touches Abraham's justifying faith or the inspiration of Sacred Scripture, the Church's teaching stands on its own authority; where commentators disagree on the hour or prophetic mode of Genesis 15, the Church leaves the question open to faithful exegetical inquiry, and those disagreements remain plainly recorded.
